\chapter{Fundamentação Teórica}\label{cap:fundamentacao}

Reúnem-se neste capítulo os conceitos necessários para entender o que este trabalho realizou e
como ele se insere no que já foi feito na área. A \autoref{sec:revisao} percorre a
literatura sobre padrões de escoamento bifásico e sobre o uso da transformada
\textit{Wavelet} em sinais de escoamento. A \autoref{sec:bases} apresenta os fundamentos
teóricos que sustentam a metodologia: o comportamento do escoamento bifásico em
microcanais e a formulação matemática das transformadas \textit{Wavelet}, incluindo a
formulação alternativa de coerência aplicada neste trabalho.

% =========================================================
\section{Revisão Bibliográfica}\label{sec:revisao}
% =========================================================

Esta seção percorre os principais trabalhos relacionados ao tema deste trabalho, com o
objetivo de situar a proposta dentro do contexto da literatura e identificar a lacuna
que justifica o trabalho.

% ---------------------------------------------------------
\subsection{Padrões de escoamento bifásico em microcanais}\label{sec:rev-padroes}
% ---------------------------------------------------------

Quando um fluido evapora dentro de um tubo de pequeno diâmetro, as duas fases não formam
uma mistura uniforme. Líquido e vapor se organizam em configurações geométricas distintas,
os chamados padrões de escoamento, que mudam conforme a quantidade de vapor presente e a
velocidade do fluido. À medida que o diâmetro diminui, os efeitos de confinamento e a
importância relativa da tensão superficial aumentam, alterando os padrões e as
condições de transição \cite{cheng2023}.

No estudo com isobutano em tubo horizontal de 1\,mm, foram identificados seis padrões:
\textit{bolha dispersa}, \textit{plug}, \textit{slug}, \textit{churn}, \textit{anular ondulado}
e \textit{anular liso}. As imagens e as respostas espectrais dos sensores forneceram
informações complementares para caracterizar esses padrões, com diferenças entre
o transdutor de pressão diferencial e o sensor óptico \cite{deoliveira2021}.

\citeonline{deoliveira2019} investigaram as séries temporais do sensor óptico durante a
ebulição convectiva de hidrocarbonetos (R-600a, R-290 e R-1270) em tubo horizontal
de 1\,mm, a 17\textdegree{}C de temperatura de saturação, com fluxo mássico entre 240 e
480\,kg/(m\textsuperscript{2}s) e fluxo de calor entre 5 e 60\,kW/m\textsuperscript{2}.
Cinco padrões foram identificados e suas frequências características analisadas por
transformada de Fourier e por métricas de não-linearidade, incluindo expoentes de Lyapunov
e retratos de fase. Os resultados indicaram comportamento caótico em todos os padrões e
frequências características distintas associadas à morfologia da interface líquido-vapor.

\citeonline{deoliveira2021} realizaram um estudo experimental mais amplo no mesmo arranjo,
com isobutano (R600a), identificando seis padrões de escoamento cujos intervalos de
ocorrência foram sistematizados em função do parâmetro adimensional $\beta$. Além do
sensor óptico (OS), analisaram-se as flutuações do transdutor de pressão diferencial
(PD), já presente no arranjo de 2019. A aquisição simultânea permitiu comparar
respostas complementares: a análise por CWT identificou frequências características
em diversas condições, mas não uma assinatura exclusiva de cada padrão em ambos
os sensores. Essas respostas são descritas na \autoref{sec:padroes}. Os parâmetros
experimentais completos estão listados na \autoref{qua:params-exp}.
Esse estudo fornece a descrição do arranjo experimental considerado neste trabalho.

\citeonline{dias2025knn} e \citeonline{dias2025cnn} desenvolveram classificadores
automáticos para os padrões do mesmo arranjo a partir de imagens de câmera de alta
velocidade. Os detalhes e resultados estão discutidos na \autoref{sec:rev-ml}.

\citeonline{schmid2022} mapearam padrões de escoamento adiabático de
CO\textsubscript{2} em tubo vertical de 8\,mm, nas direções ascendente e
descendente. O estudo mostrou a influência do título de vapor, do fluxo mássico
e da direção do escoamento sobre os padrões observados. Esses resultados
ressaltam a necessidade de considerar as condições de cada arranjo ao comparar
mapas de padrões com os de ebulição horizontal de R600a deste trabalho.

% ---------------------------------------------------------
\subsection{Análise por transformada \textit{Wavelet} em escoamentos bifásicos}\label{sec:rev-wavelet}
% ---------------------------------------------------------

A transformada \textit{Wavelet} permite identificar simultaneamente \textit{quando} e
\textit{em que frequência} determinado evento ocorre num sinal. Essa localização
complementa o espectro de Fourier global quando o conteúdo espectral varia ao longo
do registro
\cite{torrence1998}.

\citeonline{torrence1998} apresentaram um guia prático de análise \textit{Wavelet} contínua
para séries temporais, com ênfase no espectro de potência \textit{Wavelet} (WPS) e em
métodos para avaliar a significância estatística dos resultados. O artigo também introduz
o espectro cruzado (WCS) e discute as limitações do estimador clássico de coerência
\textit{Wavelet} (WCO).

Um antecedente direto da abordagem adotada neste trabalho foi desenvolvido por
\citeonline{nguyen2010}, que aplicaram a CWT a sinais de fração de vazio em escoamento
bifásico vertical ar-água em tubo de 80\,mm (laboratório KAERI, Coreia). Os autores
combinaram a energia efetiva da CWT, a razão entre o comprimento efetivo estimado
da estrutura gasosa e o diâmetro do tubo, usando a escala efetiva e a velocidade de propagação,
e a fração de vazio média.
Com esses critérios, distinguiram quatro sub-regimes: \textit{discrete bubbly},
\textit{clustered/cap bubbly}, \textit{slug} e \textit{churn}. O resultado
mostra a utilidade de combinar descritores locais da CWT com informações
físicas do escoamento \cite{nguyen2010}.

No contexto de escoamentos bifásicos em microcanais, \citeonline{deoliveira2019}
analisaram as séries temporais do sensor óptico durante a ebulição convectiva de
hidrocarbonetos no arranjo experimental da UNISINOS por meio de transformada de Fourier e
de métricas não-lineares, incluindo expoente de Lyapunov, expoente de Hurst e retratos de
fase. Os autores observaram que a morfologia da interface líquido-vapor gera assinaturas
espectrais distintas para cada padrão e que esses sinais exibem comportamento caótico
mensurável. \citeonline{deoliveira2021} aplicaram a CWT aos sinais do sensor óptico
e às flutuações de pressão diferencial do mesmo arranjo, utilizando um algoritmo
adaptado de \citeonline{torrence1998}.
Nos casos com frequências características identificáveis, o PD apresentou
frequências geralmente menores que as do OS. No padrão anular ondulado, por
exemplo, não se identificou uma frequência característica no PD, enquanto
o OS registrou oscilações associadas à interface líquido-vapor.

\citeonline{morshed2020} aplicaram decomposição por pacotes \textit{Wavelet}
com base db4 a sinais de pressão em escoamentos ar-água e ar-solução de
goma xantana. A entropia de norma evidenciou características dos regimes de
bolhas, de transição entre bolhas e bolhas alongadas e de \textit{slug}, com
limitações de resolução e resposta em frequência do sensor.
\citeonline{tan2015} combinaram sensores de
condutância e capacitância com análise por transformada \textit{Wavelet} contínua (CWT) e
demonstraram que os coeficientes de energia \textit{Wavelet} local (LWE) constituem
descritores das estruturas e transições do escoamento óleo-água, embora os autores
ressaltem que a distribuição característica obtida é específica para o arranjo e fluidos
utilizados (tubo de 50\,mm, óleo branco e água).

\citeonline{schuck2019} propuseram uma formulação alternativa para o cálculo do WPS, do
WCS e da WCO: em vez de aplicar a transformada diretamente aos sinais, aplica-se às
funções de autocorrelação e correlação cruzada. Essa mudança elimina a necessidade de
suavização espectral exigida pela definição clássica \cite{torrence1998}. Em contrapartida,
o denominador da coerência é fixo (calculado em $b = 0$, onde o WPS é máximo), enquanto o
numerador ainda varia com $b$, que agora representa o atraso (\textit{lag}) entre os
sinais, não a posição temporal absoluta. O resultado é uma representação
da relação entre dois sinais em função da escala e do
atraso. Os autores ilustraram a abordagem com três pares de realizações de sinais
sintéticos, relação sinal-ruído de 5\,dB e atraso de 500 amostras entre eventos
correspondentes, identificando o atraso e a região espectral dos pulsos.
Neste trabalho, essa formulação foi aplicada a sinais
espaciais de brilho reconstruídos a partir do panorama das imagens da câmera de alta
velocidade do arranjo descrito por \citeonline{deoliveira2021}, comparando ensaios de um mesmo
padrão de escoamento, para avaliar a repetibilidade, e ensaios de padrões distintos, para
avaliar a separação.

% ---------------------------------------------------------
\subsection{Aprendizado de máquina para classificação de padrões de escoamento}\label{sec:rev-ml}
% ---------------------------------------------------------

Nas últimas décadas, o aprendizado de máquina tornou-se uma abordagem amplamente utilizada
para identificar padrões de escoamento a partir de imagens ou sinais de sensores. Métodos
como $k$-Vizinhos Mais Próximos ($k$NN), máquinas de vetores de suporte (SVM), redes
neurais artificiais e redes neurais convolucionais (CNN) foram aplicados a diferentes
arranjos experimentais e fluidos \cite{brunton2020,yang2023}. Nos estudos do
arranjo considerado neste TCC, obtiveram-se 87\% de acurácia com $k$NN e
88{,}1\% de acurácia média com CNN \cite{dias2025knn,dias2025cnn}.
\citeonline{yang2023}, em revisão sobre visão computacional e ML aplicados a microcanais
e transferência de calor por ebulição, discutem o uso dessas técnicas na identificação
automática de padrões.

No arranjo experimental descrito por \citeonline{deoliveira2021}, \citeonline{dias2025knn}
classificaram cinco padrões de escoamento a partir de imagens da câmera de alta velocidade
com o classificador $k$NN, atingindo 87\% de acurácia. \citeonline{dias2025cnn} estenderam
esse trabalho com uma rede neural convolucional baseada na arquitetura ResNet50, alcançando
93\% no melhor \textit{fold} (média de 88{,}1\% nos dez \textit{folds}) sobre o mesmo conjunto 
de dados. No $k$NN, as imagens foram representadas por 93 atributos, incluindo
18 atributos de DWT Daubechies db2 em dois níveis; na CNN, os quadros foram
fornecidos à rede. Ambos utilizaram imagens rotuladas. Assim, a contribuição
investigada neste TCC foi a representação espacial dos panoramas e sua comparação
por espectros e coerência, complementando trabalhos que já empregavam
atributos \textit{Wavelet}.

Descritores \textit{Wavelet} de sinais unidimensionais também foram empregados
na identificação de padrões em outros arranjos experimentais.
\citeonline{sun2012} aplicaram
decomposição \textit{Wavelet} e entropia de norma sobre flutuações de pressão diferencial
em escoamento bifásico ar-água horizontal (tubo de 50\,mm). Os mapas combinaram
a entropia de norma média com a vazão mássica total ou com a fração volumétrica
de gás e obtiveram taxas globais de identificação de 98{,}0\% e 95{,}9\%,
respectivamente, para quatro padrões. Nessa avaliação, pontos sobre as fronteiras
receberam meio erro. \citeonline{tan2015} combinaram sensores de condutância e
capacitância com análise por CWT em tubo de 50\,mm e mostraram que os coeficientes de
energia \textit{Wavelet} local (LWE) caracterizam estruturas e transições do escoamento óleo-água, com a
ressalva de que a distribuição observada é específica para o arranjo e os fluidos
utilizados. Esses trabalhos motivam a investigação de descritores espectrais
dos sinais espaciais deste TCC, cuja capacidade de distinguir padrões precisou
ser examinada nos próprios dados de intensidade de imagem.

A literatura revisada aponta, portanto, dois requisitos técnicos para este trabalho:
compreender a fenomenologia dos padrões de escoamento no arranjo do LETEF e dominar
a formulação matemática da coerência \textit{Wavelet} alternativa empregada
como descritor. As próximas seções apresentam esses fundamentos.

% =========================================================
\section{Bases Teóricas}\label{sec:bases}
% =========================================================

Esta seção apresenta os conceitos teóricos que embasam a metodologia adotada, cobrindo
a fenomenologia do escoamento bifásico em microcanais e a formulação matemática das
transformadas \textit{Wavelet} utilizadas, incluindo a formulação alternativa de coerência
de \citeonline{schuck2019}.

% ---------------------------------------------------------
\subsection{Escoamento bifásico: ebulição convectiva em microcanais}\label{sec:ebolicao}
% ---------------------------------------------------------

Na ebulição convectiva em microcanais, um fluido líquido é aquecido ao percorrer um canal
de diâmetro reduzido. O aquecimento pode promover a nucleação e o crescimento de
bolhas; com a evaporação, a proporção de vapor aumenta ao longo do canal. Essa proporção é
descrita pelo título de vapor $x$, definido como a razão entre a vazão mássica de vapor e
a vazão total. Em títulos elevados, o filme líquido na parede pode secar, expondo a
parede ao vapor. A secagem do filme, denominada \textit{dryout}, e as instabilidades
do escoamento integram os fenômenos considerados na caracterização de padrões de
ebulição em microcanais \cite{cheng2023}.

O efeito relativo da tensão superficial e da gravidade pode ser expresso pelo
número de confinamento $Co$ \cite{cheng2023}:

\begin{equation}\label{eq:confinamento}
  Co=\frac{1}{D_h}\sqrt{\frac{\sigma}{g(\rho_l-\rho_v)}}.
\end{equation}

\noindent
A tensão superficial $\sigma$, a aceleração gravitacional $g$, as massas
específicas do líquido $\rho_l$ e do vapor $\rho_v$ e o diâmetro hidráulico $D_h$
determinam essa razão adimensional. Para o mesmo fluido e estado termodinâmico,
a redução do diâmetro aumenta o número de confinamento $Co$. A revisão de
\citeonline{cheng2023} reúne diferentes critérios de transição entre escalas;
neste trabalho, não se adotou um limiar universal do número de confinamento
para classificar o canal.



A transição entre padrões depende do balanço entre forças de inércia, tensão superficial
e cisalhamento. Nos ensaios em microtubo horizontal considerados na literatura deste
trabalho, foram observados padrões de bolhas, intermitentes
(\textit{plug}, \textit{slug}, \textit{churn}) e anulares
\cite{deoliveira2021,cheng2023}. As assinaturas espectrais de cada padrão nos
sinais do PD e do OS são descritas em detalhe na \autoref{sec:padroes}.

\subsubsection{Arranjo Experimental}\label{sec:arranjo}

O arranjo experimental utilizado neste trabalho foi desenvolvido no Laboratório de Estudos
Térmicos e Fluidodinâmicos (LETEF) da UNISINOS para investigar ebulição convectiva, queda
de pressão e padrões de escoamento em microcanais \cite{deoliveira2019,deoliveira2021}.
O circuito opera em malha fechada com controle independente de fluxo mássico e de fluxo de
calor, conforme esquematizado na \autoref{fig:arranjo-experimental}. O fluido refrigerante (R600a)
percorre, em sequência, uma bomba de engrenagem magnética, um \textit{subcooler}, um medidor de
vazão mássico tipo Coriolis, o pré-aquecedor (PH), a seção de teste (TS) e a seção de
visualização (VS). Após a VS, o fluido retorna ao condensador e ao vaso de líquido
antes de recircular. O reservatório de refrigerante é conectado entre esses componentes
e auxilia no controle da pressão \cite{deoliveira2021}.

\begin{figure}[htb]
  \centering
  \includegraphics[width=0.85\textwidth]{figuras/deoliveira2021-fig1-arranjo}
  \caption{Esquema do arranjo experimental do LETEF: loop fechado com bomba de engrenagem
           magnética, \textit{subcooler}, medidor Coriolis, pré-aquecedor (PH), seção de teste (TS),
           seção de visualização (VS), condensador, vaso de líquido e reservatório.}
  \label{fig:arranjo-experimental}
  \fonte{Adaptado de \citeonline{deoliveira2021}.}
\end{figure}

\paragraph{Pré-aquecedor e seção de teste.}
O pré-aquecedor (PH) e a seção de teste (TS) são tubos horizontais de aço inoxidável com
diâmetro interno de 1,0\,mm, aquecidos por efeito Joule com aplicação direta de corrente
elétrica controlada por fonte de corrente contínua (SORENSEN\texttrademark{} DCS~8-125E).
O PH tem comprimento total de 515\,mm (440\,mm
aquecidos) e estabelece a condição termodinâmica na entrada da TS. A TS tem comprimento
total de 366\,mm (265\,mm aquecidos). O aquecimento permite estabelecer diferentes
condições de ebulição na saída da TS. A jusante da TS há uma seção de
visualização (VS) de tubo de vidro com 136\,mm e mesmo diâmetro interno, onde os padrões
são registrados visualmente e o sensor óptico é posicionado \cite{deoliveira2019,deoliveira2021}.

\paragraph{Sensor óptico (OS).}
Para a análise temporal dos padrões de escoamento, foi construído um sensor óptico (OS)
que detecta a refração e a variação de luminosidade da luz na interface
líquido-vapor \cite{deoliveira2019}.
Um LED de alta luminosidade ilumina o tubo transversalmente. Um resistor dependente de luz
(LDR) no lado oposto detecta a variação de transmitância óptica causada pela passagem das
fases. A variação de resistência do LDR é convertida em tensão por um circuito com filtro
passa-baixas de quarta ordem e amplificador operacional OPA\,4227, e o sinal é adquirido
pelo módulo NI\,6009 via LabVIEW\,2013. O sensor é posicionado no início da VS, registrando
as condições da saída da TS. Por responder ao que passa pelo feixe do LED em um dado instante, 
o OS captura eventos locais e rápidos associados à morfologia da interface: a passagem de 
bolhas e frentes de vapor modifica o sinal. Em condições com muitas interfaces,
as respostas ópticas podem se sobrepor \cite{deoliveira2021}.

\paragraph{Transdutor de pressão diferencial (PD).}
A queda de pressão na TS é medida pelo transdutor ABB\texttrademark{} 600T, com frequência
de resposta de 10\,kHz \cite{deoliveira2021}. O PD mede a diferença de pressão ao
longo da seção de teste, enquanto o OS observa a transmissão óptica na seção de
visualização. Nos ensaios relatados, as frequências características do PD foram
geralmente menores que as do OS \cite{deoliveira2021}. Essa diferença observada
não equivale a uma limitação da frequência nominal de resposta do transdutor.

\paragraph{Aquisição e parâmetros experimentais.}
O OS e o PD são adquiridos \textit{simultaneamente} a 1\,kHz, com 8.000 amostras por
ensaio (8\,s por condição de operação) \cite{deoliveira2021}. Os padrões também são
registrados na seção de visualização por uma câmera de alta velocidade, cujas sequências de
quadros constituem os dados efetivamente utilizados neste trabalho. Os parâmetros dos
ensaios de referência estão resumidos na \autoref{qua:params-exp}.

\begin{quadro}[htb]
  \centering
  \caption{Parâmetros experimentais dos ensaios de \citeonline{deoliveira2021}.}
  \label{qua:params-exp}
  \begin{tabular}{ll}
    \hline
    \textbf{Parâmetro}                                        & \textbf{Valor(es)}           \\
    \hline
    Fluido de trabalho                                        & R600a (isobutano)            \\
    Temperatura de saturação [\textdegree{}C]                 & 27                           \\
    Fluxo mássico $G$ [kg/(m\textsuperscript{2}s)]           & 380, 510, 640, 760 e 890     \\
    Fluxo de calor na TS $q''$ [kW/m\textsuperscript{2}]     & 1, 5, 10, 20, 40 e 80        \\
    Frequência de aquisição [kHz]                             & 1                            \\
    Amostras por ensaio                                       & 8.000 (8\,s)                 \\
    Frequência de resposta do transdutor PD [kHz]            & 10                           \\
    \hline
  \end{tabular}
  \fonte{\citeonline{deoliveira2021}.}
\end{quadro}

As sequências de quadros da câmera de alta velocidade registradas para cada condição de
fluxo mássico e fluxo de calor constituem os dados utilizados neste trabalho: delas se
reconstrói o sinal espacial de brilho sobre o qual se calcula a coerência \textit{Wavelet}
alternativa, conforme detalhado no \autoref{cap:metodologia}.

\subsubsection{Padrões de escoamento e suas características}\label{sec:padroes}

Os padrões de escoamento observados no microtubo de 1\,mm com isobutano R600a são
descritos a seguir com base em \citeonline{deoliveira2021} e \citeonline{dias2025knn},
das configurações com bolhas dispersas às configurações anulares. As faixas adimensionais
e espectrais apresentadas correspondem às condições investigadas nesses estudos.
\citeonline{deoliveira2021} associaram cada padrão a uma faixa do parâmetro adimensional
$\beta$, definido como:

\begin{equation}\label{eq:beta}
  \beta = S_{(o\text{-}TS)} \cdot Bo \cdot 10^5,
\end{equation}

\noindent
em que $S_{(o\text{-}TS)}$ é a razão de deslizamento entre as fases na saída da seção de
teste, calculada pela razão entre as velocidades médias do vapor e do líquido.
O número de ebulição $Bo = q''_{TS}/(G\,i_{lv})$ relaciona o fluxo de calor
$q''_{TS}$ ao fluxo mássico $G$ e ao calor latente de vaporização $i_{lv}$.
O parâmetro adimensional $\beta$ combina, portanto, o deslizamento entre as
fases e o efeito do aquecimento \cite{deoliveira2021}. As faixas de $\beta$ em que cada padrão
ocorre foram estabelecidas por \citeonline{deoliveira2021} com apoio das imagens da
câmera de alta velocidade. Neste trabalho, utilizaram-se os rótulos da classificação
manual local, sem recalcular o parâmetro adimensional $\beta$ para atribuir classes.

\textbf{\textit{Bolha dispersa}} (\textit{dispersed bubble}): ocorre em baixo título de vapor e
alto fluxo mássico ($\beta < 0{,}5$)
\cite{deoliveira2021}. Pequenas bolhas esféricas ficam dispersas na fase líquida, com
diâmetro menor que o tubo. O sinal de pressão diferencial apresenta múltiplos picos entre 6 e
100\,Hz. O sensor óptico exibe energia entre 4 e 35\,Hz \cite{deoliveira2021}. Ambos os
sinais apresentam comportamento intermitente ao longo do tempo quando analisados por
\textit{Wavelet}.

\textbf{\textit{Plug}}: aparece para $0{,}5 \leq \beta < 10$, principalmente em baixos
fluxos de calor \cite{deoliveira2021}. As bolhas crescem e assumem forma cilíndrica, com
diâmetro próximo ao do tubo e comprimento maior que o diâmetro, motivo pelo qual esse padrão
também é chamado de \textit{elongated bubble flow} \cite{deoliveira2021}.
Nos ensaios relatados, também ocorreram bolhas dispersas de forma intermitente
entre os \textit{plugs} \cite{deoliveira2021}.
O sinal de pressão diferencial oscila entre 0{,}5 e 1\,Hz. O sensor óptico varia entre 2 e 16\,Hz.

\textbf{\textit{Slug}}: formado pelo coalescimento de \textit{plugs} com o aumento do título de
vapor, ocupando a faixa $10 \leq \beta < 25$ \cite{deoliveira2021}. Alterna seções de
líquido e bolsões alongados de vapor com diâmetro próximo ao do tubo
\cite{dias2025knn,deoliveira2021}. A interface é instável e libera
pequenas bolhas dispersas no líquido entre os bolsões. O estudo relatou predominância
de baixas frequências no PD. Para $G=510\,\text{kg/(m}^2\text{s)}$, os autores
relataram no OS um pico em 2{,}8\,Hz e componentes próximas de 8, 12, 23 e 32\,Hz,
atribuídas à passagem de bolhas dispersas entre os \textit{slugs}. Essas faixas
variaram com o fluxo mássico.
A natureza caótica foi confirmada com expoentes de Lyapunov entre 0{,}261
e 0{,}463 no PD e entre 0{,}413 e 4{,}362 no OS \cite{deoliveira2021}.

\textbf{\textit{Churn}}: regime de transição entre \textit{slug} e anular ($25 \leq \beta < 152$),
com título de vapor entre 0{,}35 e 0{,}62. A maior velocidade do vapor foi associada
à presença de jatos e a instabilidades na interface. A razão de deslizamento
entre as fases varia de 1{,}40 a 2{,}34. O sinal de pressão diferencial apresenta
frequências predominantes de baixa energia em 8\,Hz ou menos. No sensor óptico,
observaram-se componentes de baixa frequência e componentes mais altas, entre
15 e 28\,Hz, associadas pelos autores ao aumento do fluxo mássico \cite{deoliveira2021}.


\textbf{\textit{Anular ondulado}} (\textit{wavy annular}): observado para $152 \leq \beta < 519$.
O vapor ocupa o centro do tubo e o líquido forma um filme na parede
\cite{dias2025knn,deoliveira2021}, sobre o qual deslizam ondas frontais na interface
líquido-vapor \cite{deoliveira2021}. O sinal de pressão diferencial não apresenta
frequências características identificáveis pela CWT, enquanto o sensor óptico mantém
energia persistente entre 2 e 16\,Hz, gerada pela dinâmica dessas ondas na interface.
As flutuações de pressão apresentam expoente de Hurst abaixo de 0{,}5, indicando
comportamento anti-persistente \cite{deoliveira2021}.

\textbf{\textit{Anular liso}} (\textit{smooth annular}): ocorre para $\beta \geq 519$, em condições
extremas de baixo fluxo mássico com alto fluxo de calor \cite{deoliveira2021}. A interface
entre o filme líquido e o núcleo de vapor torna-se lisa.
\citeonline{dias2025knn} limitaram a classificação a cinco padrões
de escoamento, não incluindo o anular liso.

\begin{figure}[htb]
  \centering
  \includegraphics[width=0.90\textwidth]{figuras/dias2025knn-fig4-padroes}
  \caption{Amostras de imagens dos cinco padrões de escoamento observados no arranjo
           experimental do LETEF: (a) bolha dispersa, (b) \textit{plug}, (c) \textit{slug},
           (d) \textit{churn} e (e) anular ondulado.}
  \label{fig:padroes-escoamento}
  \fonte{Adaptado de \citeonline{dias2025knn}.}
\end{figure}

Na transição de \textit{plug} para \textit{slug}, o estudo relatou razão de
deslizamento superior a 1{,}3 e fração de vazio próxima de 0{,}65
\cite{deoliveira2021}. Esses valores caracterizam o arranjo investigado.
As diferenças entre as respostas espectrais do PD e do OS
\cite{deoliveira2021} motivam investigar representações \textit{Wavelet};
a utilidade da coerência alternativa para os panoramas foi avaliada neste TCC.



% ---------------------------------------------------------
\subsection{Transformada \textit{Wavelet}}\label{sec:wavelet}
% ---------------------------------------------------------

A transformada \textit{Wavelet} decompõe um sinal em componentes localizados ao mesmo
tempo no tempo e na frequência. A diferença em relação à transformada de Fourier é
fundamental: enquanto Fourier informa quais frequências existem no sinal, a
\textit{Wavelet} informa quando essas frequências ocorrem. Isso é feito com funções-base
chamadas \textit{Wavelets}-mãe, oscilatórias e localizadas. Ao escalar e deslocar
essa função pelo sinal, obtém-se uma representação que adapta automaticamente a resolução:
alta resolução temporal em altas frequências, boa resolução em frequência nas escalas
grandes \cite{torrence1998}. Essa adaptatividade permite analisar
sinais de escoamento bifásico, cujo comportamento muda ao longo do tempo conforme o padrão
de escoamento varia. A \autoref{qua:wavelet-fourier} compara as duas transformadas.

\begin{quadro}[htb]
  \centering
  \caption{Comparação entre a transformada de Fourier e a transformada \textit{Wavelet}.}
  \label{qua:wavelet-fourier}
  \begin{tabular}{lll}
    \hline
    \textbf{Característica}       & \textbf{Fourier global} & \textbf{Wavelet contínua}   \\
    \hline
    Localização temporal          & Não                  & Sim                            \\
    Informação principal          & Espectro global      & Conteúdo por tempo e escala    \\
    Janela de análise              & Registro completo    & Duração variável com a escala  \\
    Variação ao longo do tempo     & Não localizada       & Representada no mapa           \\
    \hline
  \end{tabular}
  \fonte{Adaptado de \citeonline{torrence1998}.}
\end{quadro}

\subsubsection{Transformada Wavelet Contínua (CWT)}\label{sec:cwt}

A CWT de um sinal $x(t)$ em relação à \textit{Wavelet}-mãe $\psi$ é definida
na forma integral por \cite{schuck2019}:

\begin{equation}\label{eq:cwt}
  W(a, b) = \frac{1}{\sqrt{a}} \int_{-\infty}^{\infty} x(t)\,
  \psi^*\!\left(\frac{t - b}{a}\right) dt,
\end{equation}

\noindent
em que $a > 0$ é o parâmetro de escala, $b$ é a posição no tempo e $(\cdot)^*$ é o
complexo conjugado. O fator $1/\sqrt{a}$ preserva a energia da função-base ao
variar a escala; não implica igualdade de amplitude dos coeficientes em escalas
distintas. Escalas grandes correspondem a baixas frequências, enquanto escalas
pequenas correspondem a altas frequências. Para a Morlet com frequência angular
central adimensional $\omega_0=6$, o período de Fourier equivalente é
$\lambda=4\pi a/(6+\sqrt{38})\approx1{,}03a$, quando a escala $a$ tem unidade de
tempo \cite{torrence1998,grinsted2004}. Se a escala é expressa em número de
amostras, inclui-se o passo de amostragem $\delta t$, obtendo-se a frequência
característica $f\approx1/(1{,}03a\delta t)$ \cite{torrence1998,grinsted2004}.
Na aplicação espacial deste trabalho, substituiu-se o passo temporal pelo passo
em pixels, e o período passou a representar um comprimento de onda espacial.

O módulo quadrático dos coeficientes da CWT forma o escalograma (\textit{scalogram}):
um mapa bidimensional da potência do sinal ao longo do tempo e das escalas.
A diferença em relação
ao espectrograma de Fourier com janela deslizante está na resolução: naquele, a janela
tem comprimento fixo $T$, de modo que eventos rápidos e lentos são analisados com a mesma
resolução temporal \cite{torrence1998}. Na CWT, a janela se adapta automaticamente: para
escalas pequenas (altas frequências), ela é estreita e captura bem eventos curtos. Para
escalas grandes (baixas frequências), ela se alarga e captura melhor a periodicidade de
eventos lentos \cite{torrence1998}. Exemplos de escalogramas dos padrões
\textit{slug} e \textit{churn} para os
dois canais do arranjo experimental (\autoref{fig:scalogram-slug} a
\autoref{fig:scalogram-churn-os}) estão apresentados na \autoref{sec:wps-wcs-wco}.

A \textit{Wavelet} de Morlet complexa utilizada na análise é
\cite{torrence1998,schuck2019}:

\begin{equation}\label{eq:morlet}
  \psi(t) = \pi^{-1/4}\, e^{-t^2/2}\, e^{j\omega_0 t},
\end{equation}

\noindent
em que $t$ é adimensional, $j=\sqrt{-1}$ e $\pi^{-1/4}$ normaliza a energia.
O parâmetro $\omega_0$ controla o equilíbrio
entre resolução temporal e espectral: valores menores aumentam a resolução temporal às
custas da resolução em frequência. A expressão sem correção de média não satisfaz
estritamente a condição de admissibilidade, mas é empregada como aproximação
prática para $\omega_0>5{,}5$ \cite{schuck2019}; neste trabalho, utilizou-se
$\omega_0=6$. Por ser complexa, a Morlet fornece coeficientes com módulo e fase,
permitindo estudar a relação entre dois sinais \cite{torrence1998}.

\subsubsection{Transformada Wavelet Discreta (DWT)}\label{sec:dwt}

A DWT é uma versão discreta e computacionalmente eficiente da análise por
\textit{Wavelet} \cite{mallat1999}. Em vez de varrer continuamente todas as escalas
e posições como na CWT, ela restringe essa varredura a uma grade diádica, o que
permite uma implementação por banco de filtros com custo computacional
$\mathcal{O}(N)$ para um sinal de $N$ amostras. O resultado é uma decomposição
em múltiplas resoluções: cada nível da decomposição captura uma faixa de frequência
diferente, do mais fino ao mais grosseiro. A DWT é apresentada aqui por completude
e por sua presença em trabalhos correlatos \cite{morshed2020,dias2025knn}. A
metodologia deste trabalho empregou exclusivamente a CWT, conforme formulação
de \citeonline{schuck2019}.

\subsubsection{Espectro de Potência, Espectro Cruzado e Coerência \textit{Wavelet}}\label{sec:wps-wcs-wco}

A partir dos coeficientes da CWT, é possível calcular três medidas espectrais que descrevem
o sinal individualmente ou a relação entre dois sinais \cite{torrence1998}.

\paragraph{Espectro de Potência Wavelet (WPS).}
O WPS mostra como a energia do sinal $x$ se distribui ao longo do tempo e das escalas
de frequência:

\begin{equation}\label{eq:wps}
  P_x(a, b) = |W_x(a, b)|^2.
\end{equation}

Regiões de alto $P_x$ indicam quando e em que frequência o sinal tem atividade intensa.
A média temporal do WPS fornece o espectro \textit{Wavelet} global, comparável a
uma versão suavizada do espectro de Fourier \cite{torrence1998}.

O transdutor de pressão diferencial (PD) mede a diferença de pressão entre as
extremidades da seção de teste, enquanto o sensor óptico (OS) registra a variação
local de luminosidade na seção de visualização. As
\autoref{fig:scalogram-slug} e \autoref{fig:scalogram-churn} mostram os escalogramas do
WPS do PD para os padrões \textit{slug} e \textit{churn}. As
\autoref{fig:scalogram-slug-os} e \autoref{fig:scalogram-churn-os} mostram os correspondentes
para o OS. Os exemplos apresentam distribuições diferentes de energia e componentes
em faixas sobrepostas. O estudo relatou, em geral, frequências características menores
no PD do que no OS \cite{deoliveira2021}.

\begin{figure}[htb]
  \centering
  \includegraphics[width=0.85\textwidth]{figuras/deoliveira2021-fig13-slug-pd}
  \caption[WPS da pressão diferencial no padrão \textit{slug}.]{Padrão \textit{slug}:
           (a) trecho do sinal de pressão diferencial para $G=760\,\text{kg/(m}^2\text{s)}$;
           escalogramas e perfis de potência média para (b) $G=510\,\text{kg/(m}^2\text{s)}$
           e (c) $G=760\,\text{kg/(m}^2\text{s)}$.}
  \label{fig:scalogram-slug}
  \fonte{Adaptado de \citeonline{deoliveira2021}.}
\end{figure}

\begin{figure}[htb]
  \centering
  \includegraphics[width=0.85\textwidth]{figuras/deoliveira2021-fig17-churn-pd}
  \caption[WPS da pressão diferencial no padrão \textit{churn}.]{Padrão \textit{churn}:
           (a) trecho do sinal de pressão diferencial para $G=510\,\text{kg/(m}^2\text{s)}$;
           escalogramas e perfis de potência média para (b) $G=510\,\text{kg/(m}^2\text{s)}$
           e (c) $G=760\,\text{kg/(m}^2\text{s)}$.}
  \label{fig:scalogram-churn}
  \fonte{Adaptado de \citeonline{deoliveira2021}.}
\end{figure}

\begin{figure}[htb]
  \centering
  \includegraphics[width=0.85\textwidth]{figuras/deoliveira2021-fig14-slug-os}
  \caption[WPS do sensor óptico no padrão \textit{slug}.]{Padrão \textit{slug}:
           (a) trecho do sinal do sensor óptico para $G=760\,\text{kg/(m}^2\text{s)}$;
           escalogramas e perfis de potência média para (b) $G=510\,\text{kg/(m}^2\text{s)}$
           e (c) $G=760\,\text{kg/(m}^2\text{s)}$.}
  \label{fig:scalogram-slug-os}
  \fonte{Adaptado de \citeonline{deoliveira2021}.}
\end{figure}

\begin{figure}[htb]
  \centering
  \includegraphics[width=0.85\textwidth]{figuras/deoliveira2021-fig18-churn-os}
  \caption[WPS do sensor óptico no padrão \textit{churn}.]{Padrão \textit{churn}:
           (a) trecho do sinal do sensor óptico; escalogramas e perfis de potência média
           para (b) $G=510\,\text{kg/(m}^2\text{s)}$ e (c) $G=760\,\text{kg/(m}^2\text{s)}$.
           Os mapas apresentam componentes de baixa frequência e componentes mais altas
           distribuídas ao longo do registro.}
  \label{fig:scalogram-churn-os}
  \fonte{Adaptado de \citeonline{deoliveira2021}.}
\end{figure}

\paragraph{Espectro Cruzado Wavelet (WCS).}
O WCS mede a relação entre dois sinais $x$ e $y$ no domínio tempo-escala:

\begin{equation}\label{eq:wcs}
  C_{xy}(a, b) = W_x(a, b) \cdot W_y^*(a, b).
\end{equation}

O WCS é o análogo \textit{Wavelet} da covariância espectral entre dois sinais: enquanto o WPS
de um sinal indica quando e em que escala ele tem energia, o WCS indica quando e em
que escala os dois sinais têm energia simultaneamente e com qual defasagem
\cite{torrence1998,grinsted2004}. Por ser um número complexo, seu módulo representa
a potência comum, ainda sem normalização pelas potências individuais, enquanto a fase descreve
sua defasagem relativa. Interpretar essa defasagem como atraso temporal requer
considerar a frequência da componente analisada \cite{grinsted2004}.

\paragraph{Coerência Wavelet (WCO).}
A WCO é a versão normalizada do WCS, com valores entre 0 e 1:

\begin{equation}\label{eq:wco-classica}
  \gamma_{xy}(a, b) = \frac{|S\bigl[a^{-1}C_{xy}(a, b)\bigr]|^2}
                           {S\bigl[a^{-1}P_x(a, b)\bigr] \cdot S\bigl[a^{-1}P_y(a, b)\bigr]},
\end{equation}

\noindent
em que $S[\cdot]$ é um operador de suavização em tempo e em escala e o fator
$a^{-1}$ pondera os espectros antes dessa suavização \cite{grinsted2004}.
A coerência clássica $\gamma_{xy}$ próxima de 1 indica forte associação na
vizinhança analisada; a coerência clássica próxima de zero indica pequena
associação linear nessa representação \cite{grinsted2004}.

\paragraph{Formulação alternativa de Schuck e Bodmann.}

Sem a suavização da \autoref{eq:wco-classica}, a coerência calculada diretamente de
um único par de transformadas é unitária onde o denominador não se anula, pois
$|W_xW_y^*|^2=|W_x|^2|W_y|^2$. Essa identidade algébrica independe da relação
física entre os sinais \cite{torrence1998}. A suavização evita essa degeneração,
mas combina informações vizinhas em tempo e escala \cite{grinsted2004}.

\citeonline{schuck2019} formularam uma alternativa baseada na CWT das funções de
correlação. Para processos estacionários, definiram a autocorrelação e a correlação
cruzada pela esperança matemática:

\begin{equation}\label{eq:autocorr}
  R_{xx}(\tau)=\operatorname{E}\{x(t+\tau)x^*(t)\},
\end{equation}

\begin{equation}\label{eq:crosscorr}
  R_{xy}(\tau)=\operatorname{E}\{x(t+\tau)y^*(t)\}.
\end{equation}

\noindent
O operador $\operatorname{E}\{\cdot\}$ representa a média estatística; o atraso $\tau$
representa a defasagem entre as amostras e o asterisco indica conjugação complexa.
A autocorrelação compara um sinal com uma versão deslocada de si mesmo; a
correlação cruzada compara dois sinais para diferentes deslocamentos. Sob a
hipótese de estacionariedade, essas funções dependem do atraso, sem depender da
origem temporal $t$ \cite{schuck2019}.

A CWT é aplicada às correlações em todo o eixo de atrasos. A aplicação deste
trabalho utilizou uma realização de cada sinal e manteve o campo cruzado
complexo da definição 32.16 de \citeonline{schuck2019}:

\begin{equation}\label{eq:wps-alt}
  \widetilde P_x(a,b)=\left|W_{\widehat R_{xx}}(a,b)\right|,
  \qquad
  \widetilde P_y(a,b)=\left|W_{\widehat R_{yy}}(a,b)\right|,
\end{equation}

\begin{equation}\label{eq:wcs-alt}
  \widetilde C_{xy}(a,b)=W_{\widehat R_{xy}}(a,b),
\end{equation}

\begin{equation}\label{eq:wco-alt}
  \widetilde\gamma_{xy}(a,b)=
  \frac{|\widetilde C_{xy}(a,b)|^2}
       {\widetilde P_x(a,0)\,\widetilde P_y(a,0)}.
\end{equation}

\noindent
Os acentos circunflexos identificam as correlações estimadas a partir do registro
finito. O campo de potência alternativa $\widetilde P_x$ é o módulo da CWT da
autocorrelação, e não seu módulo ao quadrado. O campo cruzado alternativo
$\widetilde C_{xy}$ é complexo; seu módulo descreve a intensidade da correlação
filtrada pela Morlet em cada escala. A coerência alternativa
$\widetilde\gamma_{xy}$ normaliza o módulo quadrático do campo cruzado pelas
potências dos dois sinais no atraso nulo. Essa sequência corresponde às
Equações 32.14 a 32.17 de \citeonline{schuck2019}. Há uma distinção no
original: o estimador da Equação 32.20 toma o módulo antes de calcular a
média entre realizações, enquanto a definição 32.16 conserva a quantidade
complexa. Para uma única realização, essa diferença não altera a coerência
da \autoref{eq:wco-alt}, pois seu numerador usa o módulo ao quadrado;
altera, porém, a preservação da informação de fase no campo cruzado.

A coordenada $b$ agora percorre o eixo de atrasos, não o eixo de tempo absoluto
dos sinais. Somente o denominador da coerência é avaliado em $b=0$; o numerador
permanece dependente do atraso. Dessa forma, o campo de coerência conserva
informações sobre escala e deslocamento relativo, embora não indique em qual
instante do registro original um evento ocorreu \cite{schuck2019}. Na aplicação
deste trabalho, os atrasos foram expressos em pixels do panorama reconstruído.

A ausência do operador de suavização da definição clássica não elimina a
incerteza de estimação: o método continua apoiado nas funções de correlação de
processos estacionários e nos registros utilizados para estimá-las
\cite{schuck2019}. O estimador discreto e as verificações adotadas nesta execução
são apresentados na \autoref{sec:met-wco}.

O \autoref{qua:wco-comparativo} resume as diferenças relevantes para a leitura
dos mapas.

\begin{quadro}[htb]
  \centering
  \caption{Comparação entre as formulações clássica e alternativa da coerência
           \textit{Wavelet}.}
  \label{qua:wco-comparativo}
  \begin{tabular}{p{3.1cm}p{5.1cm}p{5.1cm}}
    \hline
    \textbf{Aspecto} & \textbf{Clássica} & \textbf{Alternativa, uma realização} \\
    \hline
    Entrada da CWT & Sinais diretamente & Funções de correlação estimadas \\
    Potência & $|W_x(a,b)|^2$ & $|W_{\widehat R_{xx}}(a,b)|$ \\
    Espectro cruzado & $W_x(a,b)W_y^*(a,b)$ & $W_{\widehat R_{xy}}(a,b)$ \\
    Normalização & Espectros suavizados em tempo e escala & Potências no atraso nulo da mesma escala \\
    Coordenada $b$ & Posição no sinal & Atraso da correlação \\
    Interpretação & Associação localizada no tempo & Associação por escala e atraso \\
    \hline
  \end{tabular}
  \fonte{Elaborado pelo autor com base em \citeonline{torrence1998},
         \citeonline{grinsted2004} e \citeonline{schuck2019}.}
\end{quadro}

\citeonline{schuck2019} avaliaram a formulação com pulsos gaussianos modulados,
ruído gaussiano e razão sinal-ruído de 5\,dB, impondo atraso de 500 amostras entre
eventos. Nos exemplos apresentados, a formulação alternativa permitiu identificar
o atraso e a escala dos transientes mesmo com eventos adicionais não
correlacionados. Esse resultado fundamentou o interesse pela representação de
relações em escala e atraso, sem estabelecer que a mesma resposta ocorreria em
panoramas com artefatos de montagem. A aplicação aos ensaios de escoamento foi
avaliada pelos dados e controles descritos no \autoref{cap:metodologia}.
